\documentclass{article}
\begin{document}
\title{On the Placement of Lines}
\author{A. Writer \and B. Reader}
\date{1 January 2025}
\maketitle
\section{Start}
A typesetting system takes a description of a document and decides where every character goes. It breaks paragraphs into lines, stacks the lines into pages and places headings, notes and figures where the reader expects them. Each of these steps follows rules that can be written down precisely, and a good program follows them exactly.

Lines of a paragraph are chosen together rather than one at a time. The program looks at every way of breaking the paragraph and picks the one whose lines are spaced most evenly, avoiding hyphens on consecutive lines and very loose lines next to very tight ones.
\end{document}
